\section{Evaluation Methodology}\label{section:Evaluation_Methodology}

\makeatletter
\def\input@path{{chapters/03_methodology/06_evaluation_methodology/}}
\makeatother

\pinktodo{Add more research questions as you see fit. For example, we might add more evaluation experiments for knowledge-extraction stages other than relate}

This section defines how the empirical claims of the thesis are measured, and Chapter~\ref{chapter:Results} reports the results. The evaluation methodology is organized around the research questions: Subsection~\ref{subsec:eval-doc-extraction} supports \textbf{RQ1}, which is about the quality and correctness of the document-extraction pipeline; Subsection~\ref{subsec:eval-knowledge-extraction} supports \textbf{RQ2}, which evaluates the knowledge extraction pipeline against silver labels generated by a strong LLM, and also reports the source-provenance carry-rate to verify the provenance-preservation claims. Subsection~\ref{subsec:eval-abc} supports \textbf{RQ3}, which concerns the calibration and evaluation of the agreement-based cascading system, used in the MAP stage of the knowledge extraction. The evaluation includes voter-pool composition, and several leave-one-voter-out ablations that test each model's contribution to quality and escalation behavior. Finally, Subsection~\ref{subsec:eval-grounding-and-relation-classification} supports \textbf{RQ4}, which examines the grounding and relation-classification ablation texts. 

\pinktodo{Is the ILP-generated dataset going to be used for calibration or evaluation? Change next paragraph depending on the answer.}

Four further sections explain the evaluation methodology in more detail: how the splits for calibration and evaluation are generated, in Subsection~\ref{subsec:eval-datasets}; the offline-replay design that describes how cascade calibration works without making repeated API calls, in Subsection~\ref{subsec:eval-reprod-and-offline-replay}; how the evaluation dataset is generated using integer linear programming, explained in Subsection~\ref{subsec:eval-ilp}; and the methodological limits of the evaluation design in Subsection~\ref{subsec:eval-limitations}.

\subsection{Evaluation questions and scope}\label{subsec:eval-questions-and-scope}

Each of the four research questions maps to a different measurement procedure. Table~\ref{tab:research-questions} summarizes this mapping, and the mapping is detailed in the following sections.

\begin{table}[h]
\centering
    \begin{tabular}{@{}lll@{}}
        \toprule
        \textbf{RQ} & \textbf{Object measured} & \textbf{Defining section} \\
        \midrule
        RQ1 & Rubric-scored document-extraction faithfulness
        & Section~3.6.3 \\
        RQ2 & Knowledge-extraction strict F1 and provenance carry-rate
        & Section~3.6.4 \\
        RQ3 & Cascade cost--quality trade-off
        & Section~3.6.5 \\
        RQ4 & NLI behaviour in grounding and relation classification
        & Section~3.6.6 \\
        \bottomrule
    \end{tabular}
    \caption{Research questions and their measurement procedures.}
    \label{tab:research-questions}
\end{table}

This section does not report empirical findings; the purpose of this section is to define the evaluation design, in order to make the results in Chapter~\ref{chapter:Results} interpretable: what is measured, the metrics that were used for storing, how the datasets and the cases were sampled, and the offline-replay protocol used for cascade calibration and evaluation. 
\subsection{Evaluation datasets}\label{subsec:eval-datasets}

Section~\ref{section:Corpus} describes all evaluation datasets derived from the 977-paper ingested corpus, and Table~\ref{tab:dataset-purpose-mapping} summarizes their purpose. Three disjoint paper subsets are drawn from the corpus: the 27-paper document-extraction rubric set for RQ1; the 15-paper ILP-selected calibration cluster, containing 454 source cases and 2\,223 silver findings, for RQ2 and RQ3; and the 15-paper held-out generalization set of 285 source cases and 1\,657 silver findings, for RQ5. A separate synthetic dataset consisting of 300 claim pairs labeled as SUPPORT / CONTRADICT / UNRELATED evaluates relation classification (RQ4). Claude Opus 4.7 generates all silver findings and claim pairs.
\subsection{Document extraction evaluation}\label{subsec:eval-doc-extraction}

The document-extraction evaluation measures how accurately the pipeline extracts figures and tables from the 27-paper evaluation dataset, using a hand-labelled rubric with four dimensions: \texttt{crop}, \texttt{caption}, \texttt{footnote}, and \texttt{mask}. The crop dimension asks if the figure or table was cropped correctly -- not too big, not too small -- and is given a partial score in case of such imperfections. The caption dimension checks whether the correct caption is attached to the table or figure. The footnote dimension only applies to tables, and measures whether the table footnotes were captured within the crop. The mask dimension measures whether the artifact region to be removed is correct. Each emitted figure and table is assigned a score in $[0,1]$ for every applicable dimension; dimensions that do not apply to an item, e.g. footnotes for a figure, are marked \texttt{n/a} and are excluded from that dimension's totals. 

A per-artifact strict-F1 treats an artifact as a true positive, only when it scores perfectly on every applicable dimension. A near-miss such as \emph{crop too big} or \emph{table-in-figure} counts as an error with no partial credit on the strict-F1 score, even though it gets partial credit in some rubric dimensions. 

The per-dimension metrics are kept because these dimensions have different consequences for the downstream stages: a wrong mask, for instance, corrupts the extracted text by leaking table and figure text; a wrong caption, however, only affects the artifact metadata. Section~\ref{section:results-doc-extraction-performance} reports both the strict-F1 and per-dimension scores. 

The document-extraction configuration is selected by a staged configuration sweep, scored on the same 27-paper rubric. The sweep varies four axes: the table-detection backend (Docling-only, TATR-only, or Docling-TATR hybrid of Subsection~\ref{subsec:doc-table-and-figure-handling}); the detection confidence threshold (used by variants involving TATR); a footnote-expansion multiplier, which extends detected table crops to capture table footnotes; and the region-handling strategy (drop, pre-mask, reconstruct, or merge), together with a header-clipping test. Rather than a full joint sweep over all combinations of these axes, each stage fixes the best setting found so far and varies one axis, resulting in 31 variants in total. Each variant's strict-F1 is computed, and the best is pinned as the configuration used for production.
\subsection{Knowledge extraction evaluation}\label{subsec:eval-knowledge-extraction}

The knowledge extraction evaluation scores the pipeline's findings against the silver findings of the calibration cluster (Subsection~\ref{subsec:knowledge-extraction-calibration-set}). These silver findings are produced by sending the same extraction instructions that the voters received, to Opus 4.7, a model more capable than the cascade's best model (Sonnet 4.6). Opus 4.7 is not one of the cascade's voters; therefore, the evaluation is not a replay of the cascade's best tier. The generation process is described in Subsection~\ref{subsec:silver-labels-and-annotation-status}.

\pinktodo{run the sensitivity diagnostic.}

For each case, pipeline findings and silver findings are matched one-to-one. First, six structured fields of each finding -- (claim, subject entity, outcome entity, relation type, direction, category) -- are concatenated, and an embedding value is generated. A pipeline-silver pair is eligible to match only when its embedding cosine similarity is higher than 0.55. Two one-to-one matching strategies are available. The \textbf{greedy} matcher repeatedly takes the unmatched pair with the highest embedding similarity, until no more pairs can be matched. This procedure is simple and deterministic, but not globally optimal, since a greedily selected pair can block a better global pairing. The \textbf{optimal} matcher computes the best global assignment with the Hungarian algorithm, which mitigates the issue that the greedy matcher's selection logic would cause. The matching strategy is a swappable component; the results in Chapter~\ref{chapter:Results} are reported under the optimal (Hungarian) matcher, with the greedy matcher reported as a sensitivity diagnostic.

A matched pair is always counted as a loose true positive; it is also counted as a strict true positive, when none of the three fields -- relation type, direction, category -- disagree between findings. The finding pairs with embedding similarity scores higher than 0.55, which disagree on one or more of these fields, contribute 0.5 to the strict true-positive count. This partial contribution was added, as assigning a score of zero would be too harsh for findings that agree on the underlying claim, but differ in one categorical field.

Loose precision is defined as $n_{\mathrm{matched}} / n_{\mathrm{pipeline}}$, loose recall as $n_{\mathrm{matched}} / n_{\mathrm{silver}}$, and loose F1 as their harmonic mean. Strict precision and recall use the same denominators as their loose counterparts, but use the weighted strict true-positive counts, where relation type, direction, and category fields are also taken into account. Strict F1 is then computed as the harmonic mean of strict precision and strict recall.

Silver labels are LLM-generated and not ground truth; therefore, they inherit the limitations of the silver-generation model. The pipeline configuration is selected on the calibration set of Subsection~\ref{subsec:knowledge-extraction-calibration-set}; the knowledge extraction results reported in Chapter~\ref{chapter:Results} and discussed in Chapter~\ref{chapter:Discussion} are computed on that calibration set, and its generalization to the disjoint held-out evaluation set is reported separately as \textbf{RQ5} (Subsection~\ref{subsec:eval-heldout-generalization}).
\subsection{Agreement-based cascading evaluation}\label{subsec:eval-abc}

The agreement-based cascading evaluation compares the calibrated cascading system to the single-model baselines. The cascade configuration is selected on the calibration set of Subsection~\ref{subsec:knowledge-extraction-calibration-set}; its strict-F1 and cost are reported on that calibration cluster, and the calibrated cascade is then evaluated on the held-out evaluation set (Subsection~\ref{subsec:knowledge-extraction-eval-set}) and the results are reported separately as \textbf{RQ5} (Subsection~\ref{subsec:eval-heldout-generalization}).

The cascade configuration is selected using the offline replay procedure described in Subsection~\ref{subsec:eval-reprod-offline-replay}, in a screen-then-refine sequence: a screen over every combination of embedder (gemini, openai), agreement scorer (embedding, hybrid), and alignment strategy (greedy, Hungarian, soft\_max); a weight refinement of surviving structures (the hybrid weights and the soft\_max parameters); and a later experimental pass over the agreement threshold $\theta \in \{0.3, 0.4, \dots, 0.9\}$ and the rejection threshold $\theta_{\mathrm{reject}} \in \{0.0, 0.1, 0.2\}$. The safety gate policies of Subsection~\ref{subsec:abc-safety-gates} are then evaluated as a separate $2\times2$ ablation (single-voter policy $\times$ polarity-conflict forcing) at the pinned configuration, after the joint sweep on the other variables is complete. The best configuration values found on the calibration set are reported as empirical results in Chapter~\ref{chapter:Results}. 


% which sweeps the agreement threshold $\theta \in \{0.3, 0.4, \dots, 0.9\}$, the rejection threshold $\theta_{\mathrm{reject}} \in \{0.0, 0.1, 0.2\}$, and the different agreement scoring strategies: embedding and hybrid, described in Subsection~\ref{subsec:abc-agreement-scoring}. 

The cascade with its calibrated configuration values is compared to a set of single-model baselines. All comparisons are done under identical matcher settings on the calibration-cluster's silver findings, so that the comparisons differ only in the generated pipeline outputs. Strict F1-score of each system, as well as cost (per-chunk escalation rate and incurred cost), is reported in Chapter~\ref{chapter:Results}.

\pinktodo{Add leave-one-voter-out ablations.}

\pinktodo{Add bootstrap experiments to the code.}

\pinktodo{When the experiments are ready to run, explain what they are in here as a final paragraph.}
\subsection{Grounding and relation classification evaluation}\label{subsec:eval-grounding-and-relation-classification}

\pinktodo{Don't forget to write this section}
\subsection{Selection of evaluation dataset with integer linear programming}
\subsection{Reproducibility and offline replay}\label{subsec:eval-reprod-and-offline-replay}

A single offline batch driver produces all results reported in Chapter~\ref{chapter:Results}. The driver reads artifact files saved by the extraction pipelines. Each analysis creates one CSV and one JSON file. A global \texttt{manifest.json} file is also created, which records the status of each analysis experiment, its input and output files, and key numbers. 

The batch driver does not make any external LLM calls. Everything is read from already cached or pre-computed artifacts, including matcher embeddings, agreement-scorer embeddings, voter outputs, NLI scores, and pricing data. 

\pinktodo{Maybe add bootstrapping experiments here too, once they are added to the code and run}

The evaluation sets are frozen, and all results reported in Chapter~\ref{chapter:Results} are reproducible from the stored artifacts. A reader can re-run the offline batch and re-derive the reported results, without making any API calls, generating new embeddings, or voter outputs.
\subsection{Limitations of evaluation design}\label{subsec:eval-limitations}

The evaluation methodology has several limitations. These limitations are listed here, and their implications for interpretation are discussed in Chapter~\ref{chapter:Discussion}. 

\pinktodo{If we add 300 claim-pairs, can the limitation with relation ablation modes be solved? If yes, remove it from the limitations.}

\pinktodo{Another possible limitation can be that the 300 claim-pair synthetic dataset's related examples might be too related, contradicting examples might be too obviously contradicting. It is different than when it could be sampled from the dataset itself.}

\begin{itemize}
    \item \textbf{Silver labels are LLM-generated, not ground truth.} The precision of the silver labels is bounded by the biases of the LLM that generated them.
    \item \textbf{The document-extraction rubric set is small.} The document-extraction rubric set consists of 27 papers, manually and randomly sampled from the ingested corpus, and frozen for reproducibility. Since the paper set is small, the coverage of different figure-table extraction failure modes might be limited.
    \item \textbf{Cost-quality results depend on the selected voter models.} The reported escalation percentages and incurred costs highly depend on the chosen voter models and their settings.
    \item \textbf{The corpus is histopathology-focused.} The selected and ingested papers are all histopathology papers in English, selected from the PMC Open Access Subset. Therefore, the results might not be generalizable to other biomedical subfields, non-PMC papers, or non-English literature.
    \item \textbf{A single person annotated the document-extraction artifacts.} Inter-annotator agreement is not measured on the 27-paper document-extraction dataset, as the labels were produced by a single annotator under the rubric defined in \pinktodo{Section 3.6.3, Document-extraction evaluation}.
    \item \textbf{NLI four-mode ablations do not have per-pair accuracy labels.} Relation counts for each of the four modes are reported, but a comparison for their accuracy cannot be made, as the pairs do not have ground truth or silver labels.
\end{itemize}

Chapter~\ref{chapter:Results} reports the numerical outputs of the pipeline and the experiments. The limitations of the evaluation design are listed here, to bound the interpretations of the reported results, discussed in Chapter~\ref{chapter:Discussion}.